\documentclass[
  11pt,
  a4paper,
  oneside,
  parskip=half
]{scrreprt}

\usepackage[T1]{fontenc}
\usepackage[utf8]{inputenc}
\usepackage[ngerman]{babel}
\usepackage{lmodern}
\usepackage{microtype}
\usepackage[
  left=25mm,
  right=25mm,
  top=25mm,
  bottom=27mm,
  headheight=15pt
]{geometry}
\usepackage{xcolor}
\usepackage{graphicx}
\usepackage{booktabs}
\usepackage{tabularx}
\usepackage{longtable}
\usepackage{array}
\usepackage{enumitem}
\usepackage{listings}
\usepackage{fancyhdr}
\usepackage{titlesec}
\usepackage{pdflscape}
\usepackage{hyperref}
\usepackage{bookmark}
\usepackage{lastpage}

\definecolor{SmartDark}{HTML}{0E1319}
\definecolor{SmartPanel}{HTML}{172029}
\definecolor{SmartBlue}{HTML}{2FA8E8}
\definecolor{SmartLightBlue}{HTML}{CBEFFC}
\definecolor{SmartAmber}{HTML}{E0A03C}
\definecolor{SmartPurple}{HTML}{9D8DFB}
\definecolor{SmartGray}{HTML}{5F7688}
\definecolor{SmartLightGray}{HTML}{E8EEF5}
\definecolor{CodeBackground}{HTML}{F4F7F9}
\definecolor{CodeBorder}{HTML}{D8E1E8}
\definecolor{CodeComment}{HTML}{607D6A}
\definecolor{CodeKeyword}{HTML}{7556B8}
\definecolor{CodeString}{HTML}{9A5D21}

\hypersetup{
  colorlinks=true,
  linkcolor=SmartBlue,
  urlcolor=SmartBlue,
  citecolor=SmartBlue,
  pdftitle={Smart Homeassistant -- Technische Dokumentation},
  pdfauthor={Laneaster-ww},
  pdfsubject={Home-Assistant-Integration mit kontrollierter KI-Aktionsausführung},
  pdfkeywords={Home Assistant, KI, Conversation Agent, Policy, Ollama}
}

\pagestyle{fancy}
\fancyhf{}
\fancyhead[L]{\textcolor{SmartGray}{Smart Homeassistant}}
\fancyhead[R]{\textcolor{SmartGray}{Technische Dokumentation}}
\fancyfoot[L]{\textcolor{SmartGray}{Version 0.1.0}}
\fancyfoot[R]{\textcolor{SmartGray}{Seite \thepage\ von \pageref{LastPage}}}
\renewcommand{\headrulewidth}{0.4pt}
\renewcommand{\footrulewidth}{0pt}

\titleformat{\chapter}
  {\normalfont\huge\bfseries\color{SmartDark}}
  {\textcolor{SmartBlue}{\thechapter}}
  {0.8em}
  {}

\titleformat{\section}
  {\normalfont\Large\bfseries\color{SmartDark}}
  {\textcolor{SmartBlue}{\thesection}}
  {0.7em}
  {}

\titleformat{\subsection}
  {\normalfont\large\bfseries\color{SmartPanel}}
  {\textcolor{SmartBlue}{\thesubsection}}
  {0.6em}
  {}

\setlist[itemize]{
  leftmargin=1.6em,
  itemsep=0.25em,
  topsep=0.4em
}
\setlist[enumerate]{
  leftmargin=1.8em,
  itemsep=0.25em,
  topsep=0.4em
}

\newcolumntype{Y}{>{\raggedright\arraybackslash}X}
\newcolumntype{L}[1]{>{\raggedright\arraybackslash}p{#1}}

\lstdefinelanguage{yaml}{
  keywords={true,false,null},
  keywordstyle=\color{CodeKeyword}\bfseries,
  basicstyle=\ttfamily\small,
  sensitive=false,
  comment=[l]{\#},
  commentstyle=\color{CodeComment}\itshape,
  stringstyle=\color{CodeString},
  morestring=[b]',
  morestring=[b]"
}

\lstdefinestyle{smartcode}{
  backgroundcolor=\color{CodeBackground},
  basicstyle=\ttfamily\small,
  breaklines=true,
  breakatwhitespace=false,
  columns=fullflexible,
  frame=single,
  rulecolor=\color{CodeBorder},
  framesep=7pt,
  keepspaces=true,
  showstringspaces=false,
  tabsize=2,
  xleftmargin=0.5em,
  xrightmargin=0.5em
}

\lstset{style=smartcode}

\newcommand{\code}[1]{\texttt{#1}}
\newcommand{\projecturl}{https://github.com/Laneaster-ww/Homeassistant\_Assist}

\newenvironment{smartnote}[1]
{%
  \begin{quote}
    \color{SmartPanel}
    \textbf{\textcolor{SmartBlue}{#1}}\par\smallskip
}
{%
  \end{quote}
}

\newenvironment{smartwarning}[1]
{%
  \begin{quote}
    \color{SmartDark}
    \textbf{\textcolor{SmartAmber!80!black}{#1}}\par\smallskip
}
{%
  \end{quote}
}

\begin{document}

\begin{titlepage}
  \pagecolor{SmartDark}
  \color{SmartLightGray}

  \vspace*{18mm}

  {\Large\ttfamily\textcolor{SmartBlue}{HOME-ASSISTANT-INTEGRATION}\par}
  \vspace{9mm}

  {\fontsize{34}{40}\selectfont\bfseries Smart Homeassistant\par}
  \vspace{5mm}

  {\Large
    Kontrollierter KI-Assistent für Chat, Assist-Pipelines,
    Automationen und Dashboards
  \par}

  \vspace{18mm}

  \noindent
  \colorbox{SmartPanel}{%
    \begin{minipage}{0.89\textwidth}
      \color{SmartLightGray}
      \vspace{5mm}
      \textbf{\textcolor{SmartBlue}{TECHNISCHE DOKUMENTATION}}\par
      \vspace{3mm}
      Architektur, Installation, Konfiguration, Sicherheitsmodell,
      Betrieb und Entwicklung
      \vspace{5mm}
    \end{minipage}
  }

  \vfill

  \begin{tabular}{@{}ll@{}}
    \textcolor{SmartGray}{Projektversion} & 0.1.0 \\[2mm]
    \textcolor{SmartGray}{Dokumentstand} & 31. August 2026 \\[2mm]
    \textcolor{SmartGray}{Repository} & \href{\projecturl}{github.com/Laneaster-ww/Homeassistant\_Assist} \\[2mm]
    \textcolor{SmartGray}{Zielplattform} & Home Assistant, getestet mit 2026.6.4
  \end{tabular}

  \vspace{12mm}

  {\small\textcolor{SmartGray}{
    Das Modell schlägt vor. Die lokale Policy entscheidet.
  }\par}
\end{titlepage}

\nopagecolor
\color{SmartDark}

\chapter*{Dokumentinformationen}
\addcontentsline{toc}{chapter}{Dokumentinformationen}

\begin{tabularx}{\textwidth}{@{}L{38mm}Y@{}}
  \toprule
  \textbf{Merkmal} & \textbf{Wert} \\
  \midrule
  Dokumenttyp & Technische Projekt- und Betriebsdokumentation \\
  Projekt & Smart Homeassistant \\
  Version & 0.1.0 \\
  Status & Frühe Entwicklungsphase vor Version 1.0 \\
  Sprache & Deutsch \\
  Quellformat & \LaTeX \\
  Zielgruppe & Betreiber, Entwickler, Reviewer und technisch interessierte Anwender \\
  \bottomrule
\end{tabularx}

\section*{Zweck}

Dieses Dokument beschreibt die Funktionsweise und den sicheren Betrieb der
Home-Assistant-Integration Smart Homeassistant. Es dient als technische Referenz für
Installation, Konfiguration, Modellanbindung, Policy-Pflege, Fehlerdiagnose und
Weiterentwicklung.

\begin{smartwarning}{Lizenzstatus}
Das Repository enthält zum Zeitpunkt dieses Dokumentstands keine Lizenzdatei. Dieses
Dokument trifft keine darüber hinausgehende Aussage zu Nutzungs- oder
Weiterverteilungsrechten.
\end{smartwarning}

\tableofcontents
\clearpage

\chapter{Einleitung}

\section{Projektidee}

Smart Homeassistant erweitert Home Assistant um einen Conversation Agent, der natürliche
Sprache in Antworten oder kontrollierte Handlungsvorschläge überführt. Unterstützt werden
Textanfragen aus einem Chatfenster sowie bereits transkribierte Eingaben aus einer
Home-Assistant-Assist-Pipeline.

Das Sprachmodell ist nicht Teil der vertrauenswürdigen Ausführungsebene. Es darf lediglich
eine strukturierte Antwort erzeugen. Ob ein Vorschlag zulässig ist, entscheidet lokaler
Python-Code anhand einer expliziten Policy und zusätzlicher Schema- und
Plausibilitätsprüfungen.

\section{Ziele}

Die Integration verfolgt folgende Ziele:

\begin{itemize}
  \item verständliche Bedienung von Home Assistant mit natürlicher Sprache,
  \item Unterstützung lokaler und externer Modellanbieter,
  \item minimale Rechte durch eine Whitelist,
  \item transparente Bestätigung eingreifender oder dauerhafter Änderungen,
  \item deterministische Prüfung jeder ausführbaren Modellausgabe,
  \item lokale Dokumentationsantworten aus kontrollierten Quellen,
  \item nachvollziehbare Trennung von Vorschlag, Prüfung und Ausführung.
\end{itemize}

\section{Abgrenzung}

Smart Homeassistant ist keine allgemeine autonome Agentenplattform. Insbesondere gilt:

\begin{itemize}
  \item Das Modell kann keine beliebigen Home-Assistant-Dienste aufrufen.
  \item Das Modell erhält keinen direkten Dateisystemzugriff.
  \item Dashboards werden nicht als freies Kartenlayout durch das Modell erzeugt.
  \item Dokumentationsantworten dürfen im strikten Modus nur auf dem lokalen
        Dokumentationsschnappschuss und konfigurierten Quellen beruhen.
  \item Die Integration ersetzt keine Absicherung des Home-Assistant-Hosts.
\end{itemize}

\chapter{Funktionsumfang}

\begin{longtable}{@{}L{37mm}L{105mm}@{}}
  \toprule
  \textbf{Bereich} & \textbf{Beschreibung} \\
  \midrule
  \endfirsthead
  \toprule
  \textbf{Bereich} & \textbf{Beschreibung} \\
  \midrule
  \endhead
  Conversation Agent &
    Verwaltung von Gesprächs-ID, Verlauf, Sprache, Provider-Auswahl und offenen
    Rückfragen. \\
  Gerätesteuerung &
    Ausführung freigegebener Scripts und Services auf explizit erlaubten Entities. \\
  Automationen &
    Erzeugung eines YAML-Entwurfs, lokale Validierung und Aktivierung nach Bestätigung. \\
  Dashboards &
    Auswahl vorhandener Bereiche und Erzeugung von Views mit Home Assistants
    Bereichsstrategie. \\
  Dashboard-Löschung &
    Kontrollierte Löschung von Dashboard-Einträgen mit dem reservierten
    \code{ki-}-URL-Präfix. \\
  Dokumentationsmodus &
    Suche im lokalen Dokumentationsschnappschuss und reine Textantwort ohne
    Aktionsausführung. \\
  Modellanbieter &
    Ollama, OpenAI-kompatible APIs, Anthropic und Google Gemini. \\
  Provider-Verwaltung &
    Mehrere gespeicherte Anbieter, Standardanbieter und Auswahl pro Gespräch. \\
  Internationalisierung &
    Sichtbare, von der Integration formulierte Texte in Deutsch und Englisch. \\
  \bottomrule
\end{longtable}

\chapter{Systemarchitektur}

\section{Gesamtübersicht}

Die Architektur trennt fünf Verantwortungsbereiche:

\begin{enumerate}
  \item Eingabe über Chat oder Assist-Pipeline,
  \item Gesprächsverwaltung und Kontextaufbau,
  \item Modellzugriff über einen ausgewählten Client,
  \item deterministische Sicherheits- und Artefaktprüfung,
  \item Ausgabe als Text oder kontrollierte Änderung in Home Assistant.
\end{enumerate}

\begin{landscape}
  \thispagestyle{plain}
  \begin{figure}[p]
    \centering
    \includegraphics[width=0.96\linewidth]{../assets/architecture.png}
    \caption{Komponenten und Datenflüsse von Smart Homeassistant}
    \label{fig:architecture}
  \end{figure}
\end{landscape}

\section{Hauptfluss}

\begin{enumerate}
  \item Home Assistant übergibt eine Texteingabe an den Conversation Agent.
  \item Der Agent lädt die aktuelle Policy und den für die Sitzung gewählten Provider.
  \item Deterministische Vorfilter erkennen Antworten auf offene Bestätigungen,
        Dokumentationsfragen und bestimmte Statusfragen.
  \item Für eine normale Modellanfrage wird ein begrenzter Kontext aufgebaut.
  \item Der Provider-Adapter sendet die Anfrage an genau einen Modellendpunkt.
  \item Die Modellausgabe wird gegen das erwartete Schema validiert.
  \item Der Antworttyp bestimmt, ob Text zurückgegeben, eine Rückfrage gestellt oder ein
        lokaler Prüfpfad gestartet wird.
\end{enumerate}

\section{Antworttypen}

Das zentrale Modellausgabeformat \code{LLMOutput} unterscheidet:

\begin{table}[htbp]
  \centering
  \begin{tabularx}{\textwidth}{@{}L{42mm}Y@{}}
    \toprule
    \textbf{Typ} & \textbf{Verarbeitung} \\
    \midrule
    \code{chat} & Reine Textantwort ohne Ausführung. \\
    \code{explanation} & Erklärung vorhandener Zustände oder Konfigurationen. \\
    \code{clarification} & Rückfrage; das Gespräch bleibt geöffnet. \\
    \code{action\_plan} & Prüfung durch Policy und Action Broker. \\
    \code{automation\_draft} & Prüfung und Vorschau eines Automation-YAML-Entwurfs. \\
    \code{dashboard\_draft} & Auflösung bekannter Bereiche und Erzeugung nach Bestätigung. \\
    \code{dashboard\_delete} & Kontrollierte Löschung eines Dashboard-Eintrags mit
    reserviertem \code{ki-}-URL-Präfix. \\
    \bottomrule
  \end{tabularx}
  \caption{Arten strukturierter Modellausgaben}
\end{table}

Dokumentationsantworten nutzen einen separaten Client-Pfad und werden nicht als
\code{LLMOutput} behandelt.

\section{Kontextaufbau}

Der allgemeine Modellkontext kann folgende Informationen enthalten:

\begin{itemize}
  \item relevante aktuelle Entity-Zustände,
  \item bestehende Automationen,
  \item verfügbare Home-Assistant-Bereiche,
  \item in der Policy freigegebene Scripts, Services und Entities,
  \item löschbare Dashboard-Einträge mit dem reservierten \code{ki-}-URL-Präfix,
  \item den begrenzten Gesprächsverlauf.
\end{itemize}

Ein Relevanzfilter reduziert umfangreiche Entity- und Freigabelisten anhand der aktuellen
Anfrage. Die Prüfung bei der späteren Ausführung verwendet dennoch die vollständige
Policy.

\chapter{Installation}

\section{Voraussetzungen}

Für den mitgelieferten Entwicklungs- und Testbetrieb werden benötigt:

\begin{itemize}
  \item Docker Engine oder Docker Desktop,
  \item Docker Compose,
  \item Zugriff auf Port \code{8123},
  \item optional Zugriff auf Port \code{11434} für Ollama,
  \item alternativ Zugangsdaten für einen unterstützten externen Modellanbieter.
\end{itemize}

\section{Installation mit Docker Compose}

\begin{lstlisting}[language=bash,caption={Repository klonen und Dienste starten}]
git clone https://github.com/Laneaster-ww/Homeassistant_Assist.git
cd Homeassistant_Assist
docker compose up -d
\end{lstlisting}

Danach ist Home Assistant unter
\href{http://localhost:8123}{http://localhost:8123} erreichbar.

Die Compose-Konfiguration bindet folgende Ressourcen ein:

\begin{table}[htbp]
  \centering
  \begin{tabularx}{\textwidth}{@{}L{49mm}L{57mm}Y@{}}
    \toprule
    \textbf{Host} & \textbf{Container} & \textbf{Zweck} \\
    \midrule
    \code{./config} & \code{/config} & Home-Assistant-Laufzeitdaten und Konfiguration \\
    \code{./custom\_components/smart\_homeassistant} &
    \code{/config/custom\_components/smart\_homeassistant} &
    Integration direkt aus dem Repository \\
    benanntes Docker-Volume & \code{/root/.ollama} & lokal gespeicherte Ollama-Modelle \\
    \bottomrule
  \end{tabularx}
  \caption{Persistente und eingebundene Ressourcen}
\end{table}

\begin{smartwarning}{Vertrauliche Laufzeitdaten}
Das Verzeichnis \code{config/} kann Secrets, API-Keys, Datenbanken und interne
Home-Assistant-Daten enthalten. Es ist durch \code{.gitignore} vom Repository
ausgeschlossen und darf nicht veröffentlicht werden.
\end{smartwarning}

\section{Betriebsbefehle}

\begin{lstlisting}[language=bash,caption={Häufige Docker-Compose-Befehle}]
# Status
docker compose ps

# Home Assistant neu starten
docker compose restart homeassistant

# Protokolle verfolgen
docker compose logs -f homeassistant

# Dienste stoppen
docker compose down
\end{lstlisting}

\section{Alternative Host-Ports}

Sind die Standardports belegt, kann eine lokale und nicht versionierte
\code{docker-compose.override.yml} verwendet werden:

\begin{lstlisting}[language=yaml,caption={Beispiel für alternative Host-Ports}]
services:
  homeassistant:
    container_name: smart-homeassistant-dev
    ports: !override
      - "18123:8123"

  ollama:
    container_name: smart-homeassistant-ollama-dev
    ports: !override
      - "11435:11434"
\end{lstlisting}

\section{Manuelle Installation}

\begin{enumerate}
  \item Den Ordner \code{custom\_components/smart\_homeassistant/} in das
        Home-Assistant-Konfigurationsverzeichnis kopieren.
  \item Home Assistant vollständig neu starten.
  \item Unter \emph{Einstellungen $\rightarrow$ Geräte \& Dienste} eine neue Integration
        hinzufügen.
  \item \emph{Smart Homeassistant} auswählen und den Einrichtungsdialog abschließen.
\end{enumerate}

\begin{lstlisting}[caption={Erwartete Verzeichnisstruktur}]
config/
|-- configuration.yaml
|-- smart_homeassistant_policy.yaml
`-- custom_components/
    `-- smart_homeassistant/
        |-- __init__.py
        |-- manifest.json
        `-- ...
\end{lstlisting}

\chapter{Konfiguration}

\section{Policy-Datei}

Beim ersten Start erzeugt die Integration
\code{smart\_homeassistant\_policy.yaml}. Die Datei ist die verbindliche
Berechtigungsquelle für ausführbare Aktionen.

\begin{lstlisting}[language=yaml,caption={Reduziertes Policy-Beispiel}]
scripts:
  script.gute_nacht:
    confirmation_required: true

services:
  light.turn_on:
    confirmation_required: false
    allowed_entities:
      - light.wohnzimmer

  lock.unlock:
    confirmation_required: true
    allowed_entities:
      - lock.haustuer

automations:
  create:
    confirmation_required: true
    allowed_entities:
      - light.wohnzimmer
      - sensor.solarleistung

documentation:
  enabled: true
  strict_docs_only: true
  offline_path: custom_components/smart_homeassistant/homeassistant_docs_offline.md
\end{lstlisting}

\begin{smartwarning}{Beispielwerte ersetzen}
Die Entity-IDs im Repository sind Beispiele. Sie müssen durch vorhandene IDs der eigenen
Home-Assistant-Instanz ersetzt werden.
\end{smartwarning}

\section{Policy neu laden}

Nach einer Änderung kann die Policy ohne vollständigen Neustart über den Dienst
\code{smart\_homeassistant.reload\_policy} neu eingelesen werden.

\section{Modellanbieter}

Unterstützte Provider-Typen:

\begin{table}[htbp]
  \centering
  \begin{tabularx}{\textwidth}{@{}L{38mm}L{45mm}Y@{}}
    \toprule
    \textbf{Typ} & \textbf{Authentifizierung} & \textbf{Hinweis} \\
    \midrule
    Ollama & üblicherweise ohne API-Key & Lokaler Endpunkt, Standardadresse im
    Docker-Betrieb über \code{host.docker.internal}. \\
    OpenAI-kompatibel & API-Key abhängig vom Anbieter & Unterstützt unter anderem
    OpenAI, OpenRouter, Mistral und LM Studio. \\
    Anthropic & API-Key & Verwendet das Anthropic-Nachrichtenformat. \\
    Google Gemini & API-Key & Verwendet die Gemini-API. \\
    \bottomrule
  \end{tabularx}
  \caption{Unterstützte Modellanbieter}
\end{table}

Provider können über das Chatfenster oder die registrierten Home-Assistant-Dienste
verwaltet werden. Ein Provider kann als Standard markiert werden; zusätzlich ist eine
Auswahl pro Gespräch möglich.

\section{API-Key-Speicherung}

Providerdaten werden mit Home Assistants Storage-Helfer in \code{.storage/} gespeichert.
API-Keys liegen dort im Klartext. Sie werden bei einer Auflistung der Provider nicht an
das Frontend zurückgegeben; sichtbar ist nur, ob ein Key vorhanden ist.

\chapter{Home-Assistant-Dienste}

\begin{longtable}{@{}L{61mm}L{82mm}@{}}
  \toprule
  \textbf{Dienst} & \textbf{Zweck} \\
  \midrule
  \endfirsthead
  \toprule
  \textbf{Dienst} & \textbf{Zweck} \\
  \midrule
  \endhead
  \code{smart\_homeassistant.reload\_policy} &
    Policy ohne Home-Assistant-Neustart neu einlesen. \\
  \code{smart\_homeassistant.refresh\_documentation} &
    Konfigurierte offizielle Quellen abrufen und den lokalen
    Dokumentationsschnappschuss aktualisieren. \\
  \code{smart\_homeassistant.add\_model\_provider} &
    Modellanbieter hinzufügen. \\
  \code{smart\_homeassistant.update\_model\_provider} &
    Felder eines vorhandenen Modellanbieters aktualisieren. \\
  \code{smart\_homeassistant.remove\_model\_provider} &
    Modellanbieter entfernen. \\
  \code{smart\_homeassistant.list\_model\_providers} &
    Anbieter ohne gespeicherte API-Keys auflisten. \\
  \code{smart\_homeassistant.set\_conversation\_model} &
    Provider für ein bestimmtes Gespräch festlegen. \\
  \bottomrule
\end{longtable}

\section{Provider hinzufügen}

Wichtige Felder von \code{add\_model\_provider}:

\begin{itemize}
  \item \code{name}: frei wählbarer Anzeigename,
  \item \code{type}: \code{ollama}, \code{openai}, \code{anthropic} oder \code{gemini},
  \item \code{model}: Modellbezeichner beim Anbieter,
  \item \code{url}: optionaler oder lokaler Basisendpunkt,
  \item \code{api\_key}: falls der Provider eine Authentifizierung verlangt,
  \item \code{make\_default}: Anbieter als Standard für neue Gespräche markieren.
\end{itemize}

\section{Provider aktualisieren}

Beim Aktualisieren bleiben nicht angegebene Felder unverändert. Ein leer gelassener
API-Key ersetzt den gespeicherten Key nicht. Die Provider-ID bleibt erhalten, damit
laufende Gespräche ihre Auswahl behalten können.

\chapter{Sicherheitsmodell}

\section{Vertrauensgrenzen}

\begin{table}[htbp]
  \centering
  \begin{tabularx}{\textwidth}{@{}L{45mm}Y@{}}
    \toprule
    \textbf{Bereich} & \textbf{Einordnung} \\
    \midrule
    Nutzereingabe & Nicht vertrauenswürdig; kann mehrdeutig oder manipulativ sein. \\
    Modellausgabe & Nicht vertrauenswürdig; muss vollständig lokal validiert werden. \\
    Policy & Vom Betreiber konfigurierte Berechtigungsquelle. \\
    Broker und Artefaktprüfung & Deterministische Kontrollschicht. \\
    Modellanbieter & Konfigurierter externer oder lokaler Dienst; keine
    Ausführungsberechtigung in Home Assistant. \\
    Home-Assistant-Service-Registry & Ausführungsebene nach erfolgreicher Prüfung. \\
    \bottomrule
  \end{tabularx}
  \caption{Vertrauensgrenzen}
\end{table}

\section{Service-Aktionspläne}

Der Action Broker akzeptiert ausschließlich die Aktionen \code{run\_script} und
\code{call\_service}. Für jeden Schritt werden Service, Entity-Ziel und
Bestätigungspflicht geprüft.

Nicht ausreichend prüfbare Ziele wie \code{device\_id}, \code{area\_id},
\code{label\_id} und \code{floor\_id} werden abgelehnt, weil sie mehrere Entities
erreichen können und eine Entity-Whitelist umgehen würden.

\section{Entity-Auflösung}

Ungenaue oder falsch formatierte Entity-IDs können nur dann korrigiert werden, wenn die
Auflösung innerhalb der freigegebenen Ziele eindeutig ist. Eine solche Korrektur
erzwingt eine Bestätigung. Die Rückfrage zeigt das tatsächlich aufgelöste Ziel.

\section{Bestätigungen}

\begin{itemize}
  \item Service-Aktionspläne folgen der Bestätigungseinstellung ihrer Policy-Einträge.
  \item Korrigierte Entity-Ziele benötigen immer eine Bestätigung.
  \item Neue Automationen benötigen immer eine Bestätigung.
  \item Dashboard-Erstellung und -Löschung benötigen immer eine Bestätigung.
  \item Offene Rückfragen verfallen nach der konfigurierten Zeit.
\end{itemize}

\section{Erneute Prüfung}

Ein bestätigter Aktionsplan wird direkt vor seiner Ausführung nochmals gegen die aktuelle
Policy geprüft. Wird eine Freigabe zwischen Vorschlag und Bestätigung entfernt, findet
keine Ausführung statt.

\section{Automationsprüfung}

Automation-Entwürfe werden vor einer Bestätigung gegen folgende Kriterien geprüft:

\begin{itemize}
  \item gültiges YAML,
  \item akzeptiertes Home-Assistant-Automationsschema,
  \item erlaubte Serviceaufrufe,
  \item freigegebene Entity-Referenzen,
  \item Ausschluss nicht prüfbarer Selektoren,
  \item Ausschluss nicht erlaubter Templates,
  \item Erkennung bekannter, versehentlich kopierter Prompt-Beispiele.
\end{itemize}

\section{Dashboardprüfung}

Das Modell liefert nur Bereiche und einen Titel. Die eigentlichen Views erstellt Home
Assistant über die Bereichsstrategie. Die Integration:

\begin{itemize}
  \item löst ausschließlich bekannte Bereiche auf,
  \item erzeugt eindeutige URL-Pfade,
  \item schreibt nur in das vorgesehene Dashboard-Verzeichnis,
  \item validiert den Eintrag in \code{configuration.yaml} erneut,
  \item rollt unvollständige Änderungen bei Fehlern zurück,
  \item behandelt Dashboard-Einträge mit dem reservierten URL-Präfix \code{ki-} als
        verwaltet und löschbar.
\end{itemize}

Das Präfix \code{ki-} sollte nicht für manuell angelegte Dashboards verwendet werden,
weil solche Einträge sonst ebenfalls als Löschkandidaten erkannt werden.

\section{Transport von API-Keys}

Ist ein API-Key vorhanden, blockiert die Integration explizite
\code{http://}-Endpunkte auf nicht lokalen Hosts. Lokale Endpunkte bilden eine
Ausnahme, da sie typischerweise im eigenen Netz betrieben werden.

\begin{smartnote}{Geltungsbereich}
Diese Prüfung verhindert einen konkreten unverschlüsselten Versand an entfernte Hosts.
Sie ist keine vollständige TLS-, Zertifikats- oder Vertrauensprüfung des Anbieters.
\end{smartnote}

\section{Nicht abgedeckte Risiken}

Das Modell schützt nicht vor:

\begin{itemize}
  \item einem kompromittierten Home-Assistant-Host,
  \item absichtlich zu weit gefassten Policy-Einträgen,
  \item Schreibzugriff Dritter auf Policy, \code{.storage/} oder Integrationscode,
  \item kompromittierten Modellanbietern,
  \item Schwachstellen in Home Assistant oder aufgerufenen Integrationen.
\end{itemize}

\chapter{Dokumentationsmodus}

\section{Ablauf}

\begin{enumerate}
  \item Ein lokaler Vorfilter erkennt eine Dokumentationsfrage.
  \item Die Integration durchsucht den gespeicherten Offline-Schnappschuss.
  \item Nur passende Textabschnitte und Quellen werden an den Provider übergeben.
  \item Die Antwort wird als reine Textantwort an den Nutzer zurückgegeben.
\end{enumerate}

Wird kein passender Abschnitt gefunden, kann die Anfrage abhängig von der Erkennung als
normale Conversation-Anfrage behandelt werden. Im strikten Modus darf der
Dokumentationsclient keine unbelegte Antwort erzeugen.

\section{Aktualisierung}

Der Dienst \code{smart\_homeassistant.refresh\_documentation} lädt die in der Policy
konfigurierten offiziellen URLs erneut und ersetzt den lokalen Schnappschuss.

\begin{smartwarning}{Schreibbarer Integrationsordner}
Beim Docker-Betrieb ist der Integrationsordner absichtlich schreibbar eingebunden, da der
Aktualisierungsdienst die Offline-Dokumentation innerhalb dieses Ordners ersetzt.
\end{smartwarning}

\chapter{Betrieb und Wartung}

\section{Protokolle}

Im Docker-Betrieb:

\begin{lstlisting}[language=bash]
docker compose logs --tail 200 homeassistant
docker compose logs -f homeassistant
\end{lstlisting}

In Home Assistant können Protokolle unter
\emph{Einstellungen $\rightarrow$ System $\rightarrow$ Protokolle} eingesehen werden.
Der relevante Logger lautet:

\begin{lstlisting}
custom_components.smart_homeassistant
\end{lstlisting}

\section{Aktualisierung}

\begin{lstlisting}[language=bash]
git pull
docker compose restart homeassistant
\end{lstlisting}

Vor einer Aktualisierung sollten mindestens folgende Daten gesichert werden:

\begin{itemize}
  \item das Home-Assistant-Konfigurationsverzeichnis,
  \item \code{smart\_homeassistant\_policy.yaml},
  \item die Providerdaten unter \code{.storage/},
  \item bestehende Automationen und generierte Dashboards.
\end{itemize}

\section{Fehlerdiagnose}

\begin{longtable}{@{}L{54mm}L{89mm}@{}}
  \toprule
  \textbf{Symptom} & \textbf{Prüfung} \\
  \midrule
  \endfirsthead
  \toprule
  \textbf{Symptom} & \textbf{Prüfung} \\
  \midrule
  \endhead
  Integration nicht auffindbar &
    Verzeichnisstruktur prüfen und Home Assistant vollständig neu starten. \\
  Aktion wird abgelehnt &
    Service, Script und Entity in der Policy kontrollieren; anschließend Policy neu
    laden. \\
  Anbieter nicht erreichbar &
    Basis-URL, Modellname, Netzwerkzugriff und Containerstatus prüfen. \\
  Unsicherer Endpunkt &
    Bei einem entfernten Host HTTPS verwenden oder den API-Key entfernen. \\
  Automationsentwurf wird verworfen &
    Protokoll auf Schema-, Service-, Entity- oder Templatefehler prüfen. \\
  Dashboard erscheint nicht &
    Von der Integration angebotenen Home-Assistant-Neustart durchführen. \\
  Dokumentationsantwort fehlt &
    Dokumentationsmodus und Offline-Pfad in der Policy kontrollieren; Schnappschuss
    aktualisieren. \\
  \bottomrule
\end{longtable}

\chapter{Entwicklerdokumentation}

\section{Repository-Struktur}

\begin{lstlisting}
Homeassistant_Assist/
|-- custom_components/
|   `-- smart_homeassistant/
|       |-- __init__.py
|       |-- conversation.py
|       |-- policy.py
|       |-- broker.py
|       |-- automations.py
|       |-- dashboard.py
|       |-- providers.py
|       |-- llm_clients.py
|       |-- ollama_client.py
|       |-- docs_knowledge.py
|       |-- config_flow.py
|       |-- sensor.py
|       |-- services.yaml
|       |-- strings.json
|       |-- frontend/
|       `-- translations/
|-- docs/
|-- examples/
|-- docker-compose.yml
`-- README.md
\end{lstlisting}

\section{Modulverantwortung}

\begin{longtable}{@{}L{48mm}L{95mm}@{}}
  \toprule
  \textbf{Modul} & \textbf{Verantwortung} \\
  \midrule
  \endfirsthead
  \toprule
  \textbf{Modul} & \textbf{Verantwortung} \\
  \midrule
  \endhead
  \code{conversation.py} &
    Gesprächsfluss, Sitzungen, Vorfilter, Rückfragen und Antwortverzweigung. \\
  \code{config\_flow.py} &
    Einrichtungs- und Optionen-Dialog, Einzelinstanz-Regel sowie Bestätigungs-TTL. \\
  \code{policy.py} &
    Laden, Validieren und Abfragen der lokalen Whitelist. \\
  \code{broker.py} &
    Prüfung und Ausführung einzelner Schritte eines Service-Aktionsplans. \\
  \code{automations.py} &
    Normalisierung, Schema- und Policyprüfung sowie Aktivierung von Automationen. \\
  \code{dashboard.py} &
    Auflösung von Bereichen, Erzeugung, Konfigurationsprüfung und Löschung von
    Dashboards. \\
  \code{providers.py} &
    Persistenz, Standardauswahl und sichere Frontend-Repräsentation der Provider. \\
  \code{llm\_clients.py} &
    Clients für OpenAI-kompatible, Anthropic- und Gemini-Endpunkte. \\
  \code{ollama\_client.py} &
    Ollama-Zugriff, Antwortschema, Prompt und gemeinsame Hilfsfunktionen. \\
  \code{docs\_knowledge.py} &
    Klassifikation, Suche und Aktualisierung der Offline-Dokumentation. \\
  \code{entity\_filter.py} &
    Relevanzauswahl umfangreicher Entity- und Policy-Listen. \\
  \code{sensor.py} &
    Nicht pollender Sensor für die zuletzt verarbeitete Aktion und den kurzen
    Aktionsverlauf. \\
  \code{frontend.py} &
    Registrierung der statischen Benutzeroberfläche. \\
  \code{services.yaml} &
    Metadaten und UI-Beschreibungen der registrierten Home-Assistant-Dienste. \\
  \code{strings.json}, \code{translations/} &
    Texte für Einrichtung, Optionen, Dienste und Übersetzungen. \\
  \code{text\_format.py}, \code{i18n.py} &
    Normalisierung von Modelltexten und sprachabhängige Laufzeitmeldungen. \\
  \code{files.py}, \code{locks.py} &
    Gemeinsame Dateioperationen und Synchronisation konkurrierender Änderungen. \\
  \bottomrule
\end{longtable}

\section{Lokaler Entwicklungsablauf}

\begin{lstlisting}[language=bash]
# Dienste starten
docker compose up -d

# Nach Python-Aenderungen
docker compose restart homeassistant

# Syntaxpruefung
python -m compileall custom_components/smart_homeassistant

# Protokolle
docker compose logs -f homeassistant
\end{lstlisting}

\section{Regeln für Änderungen}

\begin{itemize}
  \item Sicherheitsentscheidungen müssen deterministisch bleiben.
  \item Neue Aktionen benötigen eine explizite Policy-Begrenzung.
  \item Modellausgaben dürfen nicht ungeprüft an Home Assistant oder das Dateisystem
        weitergereicht werden.
  \item Sichtbare Texte müssen bei Bedarf in beiden Übersetzungsdateien ergänzt werden.
  \item Änderungen an Konfiguration oder Verhalten müssen in README und Dokumentation
        nachvollzogen werden.
  \item Beispiele dürfen keine echten Secrets oder privaten Installationsdaten enthalten.
\end{itemize}

\chapter{Datenschutz und Betriebsdaten}

\section{Lokal gespeicherte Daten}

\begin{table}[htbp]
  \centering
  \begin{tabularx}{\textwidth}{@{}L{49mm}Y@{}}
    \toprule
    \textbf{Daten} & \textbf{Speicherort} \\
    \midrule
    Providerkonfiguration und API-Keys & Home Assistant \code{.storage/} \\
    Sicherheits-Policy & \code{smart\_homeassistant\_policy.yaml} \\
    Automationen & Home-Assistant-Automationskonfiguration \\
    generierte Dashboards & \code{smart\_homeassistant\_dashboards/} und
    \code{configuration.yaml} \\
    Offline-Dokumentation &
    \code{custom\_components/smart\_homeassistant/homeassistant\_docs\_offline.md} \\
    Ollama-Modelle & benanntes Docker-Volume \\
    \bottomrule
  \end{tabularx}
  \caption{Persistente Daten}
\end{table}

\section{Externe Übertragung}

Bei Nutzung eines Cloud-Providers werden Modellkontext und Nutzereingabe an den
konfigurierten Endpunkt übertragen. Dazu können je nach Anfrage Entity-Namen,
Entity-Zustände, Automationsnamen, Bereiche und Gesprächsverlauf gehören.

Betreiber müssen selbst bewerten, welche Daten an den gewählten Provider übertragen
werden dürfen. Für maximale lokale Verarbeitung kann Ollama verwendet werden.

\chapter{Zusammenfassung}

Smart Homeassistant verbindet natürliche Sprache mit einer lokal kontrollierten
Home-Assistant-Ausführung. Die zentrale Eigenschaft ist nicht die Modellanbindung,
sondern die Trennung zwischen unverbindlichem Modellvorschlag und deterministischer
Freigabe.

Für einen sicheren Betrieb sind insbesondere wichtig:

\begin{enumerate}
  \item die Policy eng und installationsspezifisch pflegen,
  \item dauerhafte Änderungen vor der Bestätigung prüfen,
  \item das Home-Assistant-Konfigurationsverzeichnis schützen,
  \item externe Provider und die übertragenen Kontextdaten bewusst auswählen,
  \item Sicherungen und Protokollkontrolle in den regulären Betrieb aufnehmen.
\end{enumerate}

\appendix

\chapter{Kurzreferenz}

\section{Wichtige URLs und Pfade}

\begin{tabularx}{\textwidth}{@{}L{58mm}Y@{}}
  \toprule
  \textbf{Eintrag} & \textbf{Wert} \\
  \midrule
  Repository & \href{\projecturl}{github.com/Laneaster-ww/Homeassistant\_Assist} \\
  Home Assistant im Docker-Standardbetrieb & \url{http://localhost:8123} \\
  Ollama im Docker-Standardbetrieb & \url{http://localhost:11434} \\
  Integrationsdomain & \code{smart\_homeassistant} \\
  Policy-Datei & \code{smart\_homeassistant\_policy.yaml} \\
  Integrationsverzeichnis & \code{custom\_components/smart\_homeassistant/} \\
  Dashboard-Verzeichnis & \code{smart\_homeassistant\_dashboards/} \\
  \bottomrule
\end{tabularx}

\section{Provider-Typen}

\begin{center}
  \code{ollama}
  \quad
  \code{openai}
  \quad
  \code{anthropic}
  \quad
  \code{gemini}
\end{center}

\section{Weiterführende Repository-Dokumente}

\begin{itemize}
  \item \code{README.md}
  \item \code{docs/INSTALLATION.md}
  \item \code{docs/ARCHITECTURE.md}
  \item \code{docs/SECURITY\_MODEL.md}
  \item \code{CONTRIBUTING.md}
  \item \code{SECURITY.md}
  \item \code{CHANGELOG.md}
\end{itemize}

\end{document}
